\section{Ranks and the complete finite weight records}

The two prime-rank lemmas in the source, and its surrounding assertion
about other characteristics, can be handled together. For the fixed
matrices the full answer is
\[
\begin{array}{c|l|c}
n&\text{exceptional primes and ranks}&\text{rank at every other prime}\\\hline
105&2:104,\quad3:20&105\\
120&3:21,\quad7:119,\quad11:99&120\\
176&3:21,\quad5:175,\quad7:175,\quad17:155&176\\
253&2:252,\quad3:23,\quad5:230,\quad7:252&253.
\end{array}                                                     \tag{3}
\]
Indeed (2) determines the integer determinant as the product of the
listed eigenvalues with their multiplicities. Outside its prime divisors
reduction is invertible. For every prime divisor, the complete supplied
elimination checks the rank of the actual matrix. Thus the finite table
covers every exceptional characteristic; no sample of small primes is
being extrapolated. For the ternary entries, (1) also gives a written
incidence proof for $n=120,176,253$, while the $105$ entry is certified
by a complete row reduction.

This proves Lemmas~1 and~2, including the full ranks at $p=11,23$ in
Lemma~2. The paragraph preceding Lemma~1 says that the only nonternary
exception to full rank or codimension one occurs for $(n,p)=(253,5)$.
The additional entries $(120,11)$ and $(176,17)$ in (3) correct that
sentence. They do not change any of the ternary principal codes.

The distances used below come from the previously completed full
coefficient enumerations in the exact row-reduced bases, together with
the two separately completed enumerations for $K$ and $N$ at length $120$.
These are mathematical input certificates. For a basis of dimension $k$,
the enumerator visits each of the $3^k$ coefficient vectors once, and
basis independence makes their images all distinct codewords. The
complete distributions are retained, not merely their minimum entries.
The current checker compares the reconstructed bases with those exact
bases. No large enumeration is repeated during assembly.

The common MacWilliams theorem \cite[Theorem~7.1]{U} then gives every dual
coefficient over the integers. In particular the resulting parameters
and minimum multiplicities are
\[
\begin{array}{c|r|r|r|r|r|r}
\text{code}&n&k&d&A_d&d^\perp&A_{d^\perp}^\perp\\\hline
C_{105}&105&20&33&770&6&15680\\
D_{105}&105&19&36&3360&5&84\\
C_{120}&120&21&30&112&6&33600\\
S&120&20&30&112&6&43680\\
K&120&15&48&1260&4&2100\\
N&120&16&48&1260&4&1260\\
C_{176}&176&21&56&924&6&184800\\
C_{253}&253&23&77&46&6&850080\\
D_{253}&253&22&97&3542&6&850080.
\end{array}                                                     \tag{4}
\]
The dual dimension is $n-k$ in each row. Nonnegative integer coefficients
and the dual total are checked, but those consistency checks alone do
not certify an arbitrary proposed input distribution. The complete
coefficient enumeration, its matrix basis, and its method are all retained.
The two $105$ distributions agree entry by entry with source Tables~1--2;
those full distributions are reproduced in the data table accompanying
this application. The two further $20$-dimensional codes of Remark~2 use
a smaller, justified search, described below.
